\begin{thebibliography}{99}
\bibitem{imielinski} Tomasz Imieli\'nski and Witold Lipski. 1984. Incomplete information in relational databases. \textit{Journal of the ACM} 31(4), 761--791. \url{https://doi.org/10.1145/1634.1886}
\bibitem{libkin2014} Leonid Libkin. 2014. Incomplete data. In \textit{PODS 2014}, 1--13. \url{https://doi.org/10.1145/2594538.2594561}
\bibitem{ullmann} Julian R. Ullmann. 1976. An algorithm for subgraph isomorphism. \textit{Journal of the ACM} 23(1), 31--42. \url{https://doi.org/10.1145/321921.321925}
\bibitem{cordella} Luigi P. Cordella, Pasquale Foggia, Carlo Sansone, and Mario Vento. 2004. A (sub)graph isomorphism algorithm for matching large graphs. \textit{IEEE Transactions on Pattern Analysis and Machine Intelligence} 26(10), 1367--1372. \url{https://doi.org/10.1109/TPAMI.2004.75}
\bibitem{vf3} Vincenzo Carletti, Pasquale Foggia, Alessia Saggese, and Mario Vento. 2017. Introducing VF3: A new algorithm for subgraph isomorphism. In \textit{GbRPR 2017}, 128--139. \url{https://doi.org/10.1007/978-3-319-58961-9_12}
\bibitem{turboiso} Wook-Shin Han, Jinsoo Lee, and Jeong-Hoon Lee. 2013. TurboISO: Towards ultrafast and robust subgraph isomorphism search in large graphs. In \textit{SIGMOD 2013}, 337--348. \url{https://doi.org/10.1145/2463676.2465300}
\bibitem{taming} Haichuan Shang, Ying Zhang, Xuemin Lin, and Jeffrey Xu Yu. 2008. Taming verification hardness: An efficient algorithm for testing subgraph isomorphism. \textit{PVLDB} 1(1), 364--375. \url{https://doi.org/10.14778/1453856.1453899}
\bibitem{daf} Myoungji Han, Hyunjoon Kim, Geonmo Gu, Kunsoo Park, and Wook-Shin Han. 2019. Efficient Subgraph Matching: Harmonizing Dynamic Programming, Adaptive Matching Order, and Failing Set Together. In \textit{SIGMOD 2019}, 1429--1446. \url{https://doi.org/10.1145/3299869.3319880}
\bibitem{bice} Yunyoung Choi, Kunsoo Park, and Hyunjoon Kim. 2023. BICE: Exploring Compact Search Space by Using Bipartite Matching and Cell-Wide Verification. \textit{PVLDB} 16(9), 2186--2198. \url{https://doi.org/10.14778/3598581.3598591}
\bibitem{gup} Junya Arai, Yasuhiro Fujiwara, and Makoto Onizuka. 2023. GuP: Fast Subgraph Matching by Guard-based Pruning. \textit{Proceedings of the ACM on Management of Data} 1(2), 1--26. \url{https://doi.org/10.1145/3589312}
\bibitem{asdmatch} Qiyan Li, Jeffrey Xu Yu, and Zongyan He. 2025. Subgraph Matching: A New Decomposition Based Approach. \textit{PVLDB} 18(11), 4282--4294. \url{https://doi.org/10.14778/3749646.3749693}
\bibitem{cemr} Linglin Yang et al. 2026. CEMR: An Effective Subgraph Matching Algorithm with Redundant Extension Elimination. \textit{PVLDB} 19(7), 1572--1585. \url{https://doi.org/10.14778/3801059.3801070}
\bibitem{sunluo} Shixuan Sun and Qiong Luo. 2020. In-memory subgraph matching: An in-depth study. In \textit{SIGMOD 2020}, 1083--1098. \url{https://doi.org/10.1145/3318464.3380581}
\bibitem{rapidmatch} Shixuan Sun, Xibo Sun, Yulin Che, Qiong Luo, and Bingsheng He. 2020. RapidMatch: A Holistic Approach to Subgraph Query Processing. \textit{PVLDB} 14(2), 176--188. \url{https://doi.org/10.14778/3425879.3425888}
\bibitem{lad2010} Christine Solnon. 2010. AllDifferent-based filtering for subgraph isomorphism. \textit{Artificial Intelligence} 174, 12--13, 850--864. \url{https://doi.org/10.1016/j.artint.2010.05.002}
\bibitem{lad} Christine Solnon. 2026. LAD2025, A constraint-based solver for the subgraph isomorphism problem. \textit{Artificial Intelligence} 352, 104474. \url{https://doi.org/10.1016/j.artint.2025.104474}
\bibitem{glasgow} Ciaran McCreesh, Patrick Prosser, and James Trimble. 2020. The Glasgow Subgraph Solver: Using constraint programming to tackle hard subgraph isomorphism problem variants. In \textit{ICGT}, 316--324. \url{https://doi.org/10.1007/978-3-030-51372-6_19}
\bibitem{side_constraints} \"{O}zg\"{u}r Akg\"{u}n, Jessica Enright, Christopher Jefferson, Ciaran McCreesh, Patrick Prosser, and Steffen Zschaler. 2021. Finding subgraphs with side constraints. In \textit{CPAIOR}, LNCS 12735, 348--364. \url{https://doi.org/10.1007/978-3-030-78230-6_22}
\bibitem{certified} Stephan Gocht, Ciaran McCreesh, and Jakob Nordstr\"om. 2020. Subgraph isomorphism meets cutting planes: Solving with certified solutions. In \textit{IJCAI 2020}, 1134--1140. \url{https://doi.org/10.24963/ijcai.2020/158}
\bibitem{e2everify} Stephan Gocht, Ciaran McCreesh, Magnus O. Myreen, Jakob Nordstr\"om, Andy Oertel, and Yong Kiam Tan. 2024. End-to-End Verification for Subgraph Solving. In \textit{AAAI 2024}, 8038--8047. \url{https://doi.org/10.1609/aaai.v38i8.28642}
\bibitem{auditablecp} Stephan Gocht, Ciaran McCreesh, and Jakob Nordstr\"{o}m. 2022. An auditable constraint programming solver. In \textit{CP}, LIPIcs 235, 25:1--25:18. \url{https://doi.org/10.4230/LIPIcs.CP.2022.25}
\bibitem{kotthoff} Lars Kotthoff, Ciaran McCreesh, and Christine Solnon. 2016. Portfolios of Subgraph Isomorphism Algorithms. In \textit{LION 10}, LNCS 10079, 107--122. \url{https://doi.org/10.1007/978-3-319-50349-3_8}
\bibitem{mccreesh_hard} Ciaran McCreesh, Patrick Prosser, Christine Solnon, and James Trimble. 2018. When subgraph isomorphism is really hard, and why this matters for graph databases. \textit{Journal of Artificial Intelligence Research} 61, 723--759. \url{https://doi.org/10.1613/jair.5768}
\bibitem{libkin2011} Leonid Libkin. 2011. Incomplete information and certain answers in general data models. In \textit{PODS 2011}, 59--70. \url{https://doi.org/10.1145/1989284.1989294}
\bibitem{barcelo} Pablo Barcel\'o, Leonid Libkin, and Juan L. Reutter. 2014. Querying Regular Graph Patterns. \textit{Journal of the ACM} 61(1), 1--54. \url{https://doi.org/10.1145/2559905}
\bibitem{arenas} Marcelo Arenas, Leopoldo Bertossi, and Jan Chomicki. 1999. Consistent query answers in inconsistent databases. In \textit{PODS 1999}, 68--79. \url{https://doi.org/10.1145/303976.303983}
\bibitem{equip} Phokion G. Kolaitis, Enela Pema, and Wang-Chiew Tan. 2013. Efficient querying of inconsistent databases with binary integer programming. \textit{Proc. VLDB Endow.} 6, 6, 397--408. \url{https://doi.org/10.14778/2536336.2536337}
\bibitem{cavsat} Akhil A. Dixit and Phokion G. Kolaitis. 2019. A SAT-based system for consistent query answering. In \textit{SAT}, 117--135. \url{https://doi.org/10.1007/978-3-030-24258-9_8}
\bibitem{certain_answers} Giovanni Amendola and Leonid Libkin. 2018. Explainable Certain Answers. In \textit{Proceedings of IJCAI 2018}, 1683--1690. \url{https://doi.org/10.24963/ijcai.2018/233}
\bibitem{backbones} Mikol\'a\v{s} Janota, In\^es Lynce, and Joao Marques-Silva. 2015. Algorithms for computing backbones of propositional formulae. \textit{AI Communications} 28(2), 161--177. \url{https://doi.org/10.3233/AIC-140640}
\bibitem{cadiback} Armin Biere, Nils Froleyks, and Wenxi Wang. 2023. CadiBack: Extracting backbones with CaDiCaL. In \textit{SAT}, LIPIcs 271, 3:1--3:12. \url{https://doi.org/10.4230/LIPIcs.SAT.2023.3}
\bibitem{cautious} Mario Alviano, Carmine Dodaro, and Francesco Ricca. 2014. Anytime Computation of Cautious Consequences in Answer Set Programming. \textit{Theory and Practice of Logic Programming} 14(4--5), 755--770. \url{https://doi.org/10.1017/S1471068414000325}
\bibitem{clingo} Martin Gebser, Roland Kaminski, Benjamin Kaufmann, and Torsten Schaub. 2019. Multi-shot ASP solving with clingo. \textit{Theory and Practice of Logic Programming} 19(1), 27--82. \url{https://doi.org/10.1017/S1471068418000054}
\bibitem{karp} Richard M. Karp. 1972. Reducibility among combinatorial problems. In Raymond E. Miller and James W. Thatcher (eds.), \textit{Complexity of Computer Computations}. Plenum, 85--103. \url{https://doi.org/10.1007/978-1-4684-2001-2_9}
\bibitem{marx} D\'aniel Marx and Micha\l Pilipczuk. 2014. Everything you always wanted to know about the parameterized complexity of subgraph isomorphism (but were afraid to ask). In \textit{STACS 2014}, 542--553. \url{https://doi.org/10.4230/LIPIcs.STACS.2014.542}
\bibitem{shacl} Julien Corman, Juan L. Reutter, and Ognjen Savkovi\'{c}. 2018. Semantics and validation of recursive SHACL. In \textit{ISWC}, LNCS 11136, 318--336. \url{https://doi.org/10.1007/978-3-030-00671-6_19}
\bibitem{pasula} Hanna Pasula, Bhaskara Marthi, Brian Milch, Stuart Russell, and Ilya Shpitser. 2002. Identity uncertainty and citation matching. In \textit{NIPS}, 1401--1408.
\bibitem{moustafa} Walaa Eldin Moustafa, Angelika Kimmig, Amol Deshpande, and Lise Getoor. 2014. Subgraph pattern matching over uncertain graphs with identity linkage uncertainty. In \textit{IEEE ICDE 2014}, 904--915. \url{https://doi.org/10.1109/ICDE.2014.6816710}
\bibitem{reiter} Raymond Reiter. 1987. A theory of diagnosis from first principles. \textit{Artificial Intelligence} 32(1), 57--95. \url{https://doi.org/10.1016/0004-3702(87)90062-2}
\bibitem{peng} Roger D. Peng. 2011. Reproducible research in computational science. \textit{Science} 334(6060), 1226--1227. \url{https://doi.org/10.1126/science.1213847}
\bibitem{gundersen} Odd Erik Gundersen and Sigbj\o rn Kjensmo. 2018. State of the art: Reproducibility in artificial intelligence. In \textit{AAAI 2018}. \url{https://doi.org/10.1609/aaai.v32i1.11503}
\bibitem{pineau} Joelle Pineau et al. 2021. Improving reproducibility in machine learning research: A report from the NeurIPS 2019 reproducibility program. \textit{Journal of Machine Learning Research} 22(164), 1--20.
\bibitem{datasheets} Timnit Gebru, Jamie Morgenstern, Briana Vecchione, Jennifer Wortman Vaughan, Hanna Wallach, Hal Daum\'e III, and Kate Crawford. 2021. Datasheets for datasets. \textit{Communications of the ACM} 64(12), 86--92. \url{https://doi.org/10.1145/3458723}
\bibitem{networkx} Aric A. Hagberg, Daniel A. Schult, and Pieter J. Swart. 2008. Exploring Network Structure, Dynamics, and Function using NetworkX. In \textit{Proceedings of the 7th Python in Science Conference (SciPy 2008)}, 11--15.
\bibitem{hprd} T. S. Keshava Prasad et al. 2009. Human Protein Reference Database---2009 update. \textit{Nucleic Acids Research} 37(suppl\_1), D767--D772. \url{https://doi.org/10.1093/nar/gkn892}
\bibitem{cora} Prithviraj Sen, Galileo Namata, Mustafa Bilgic, Lise Getoor, Brian Gallagher, and Tina Eliassi-Rad. 2008. Collective classification in network data. \textit{AI Magazine} 29, 3, 93--106. \url{https://doi.org/10.1609/aimag.v29i3.2157}
\bibitem{pubmed} Galileo Namata, Ben London, Lise Getoor, and Bert Huang. 2012. Query-driven active surveying for collective classification. In \textit{Workshop on Mining and Learning with Graphs (MLG)}.
\bibitem{planetoid} Zhilin Yang, William W. Cohen, and Ruslan Salakhutdinov. 2016. Revisiting semi-supervised learning with graph embeddings. In \textit{ICML}, PMLR 48, 40--48.
\bibitem{cadical} Armin Biere, Tobias Faller, Katalin Fazekas, Mathias Fleury, Nils Froleyks, and Florian Pollitt. 2024. CaDiCaL 2.0. In \textit{CAV}, LNCS 14681, 133--152. \url{https://doi.org/10.1007/978-3-031-65627-9_7}
\bibitem{drattrim} Nathan Wetzler, Marijn J.~H. Heule, and Warren A. Hunt Jr. 2014. DRAT-trim: Efficient checking and trimming using expressive clausal proofs. In \textit{SAT}, 422--429. \url{https://doi.org/10.1007/978-3-319-09284-3_31}
\end{thebibliography}
